\subsection{Everett's finite parity coordinates and density-one descent}
\label{subsec:Everett-1977}

Everett works on \(\Nzero\) with the map
\begin{equation}
\label{eq:Everett-f}
 f_{\mathrm{Ev}}(m)=
 \begin{cases}
  m/2,&m\text{ even},\\
  (3m+1)/2,&m\text{ odd}.
 \end{cases}
\end{equation}
The subscript is editorial: the source writes \(f\).  Thus
\(f_{\mathrm{Ev}}=T|_{\Nzero}\), exactly; it is neither the unshortened map
\(U\) nor Crandall's odd first-return map.  Put
\begin{equation}
\label{eq:Everett-parity-coordinates}
 m_n=f_{\mathrm{Ev}}^n(m),
 \qquad
 x_n=m_n\bmod2\in\{0,1\}.
\end{equation}
The source word is chronological.  When encoded as an integer below, \(x_0\)
is the least significant digit \cite[pp.~42--43]{Everett1977}.

\begin{theorem}[Everett, Theorem 1: finite parity fibres]
\label{thm:Everett-finite-fibres}
Let \(N\ge1\), let
\(w=(x_0,\ldots,x_{N-1})\in\{0,1\}^N\), and put
\[
 X=X(w)=\sum_{n=0}^{N-1}x_n.
\]
There are unique integers \(a(w),b(w)\) satisfying
\[
 0\le a(w)<2^N,
 \qquad
 0\le b(w)<3^X,
\]
such that the complete nonnegative fibre of \(w\) and its \(N\)-th iterate
are
\begin{equation}
\label{eq:Everett-affine-fibre}
 m=a(w)+2^NQ,
 \qquad
 m_N=b(w)+3^XQ,
 \qquad Q\in\Nzero.
\end{equation}
Consequently every length-\(N\) binary word occurs at one and only one
representative \(0\le m<2^N\).
\end{theorem}

\begin{proof}
This is the source induction with its implicit quotient domain made explicit.
For \(N=1\), \(x_0=0\) gives
\[
 m=2Q,\qquad m_1=Q,\qquad (a,b,X)=(0,0,0),
\]
whereas \(x_0=1\) gives
\[
 m=1+2Q,\qquad m_1=2+3Q,\qquad (a,b,X)=(1,2,1).
\]

Assume \eqref{eq:Everett-affine-fibre} at length \(N\), and abbreviate its
coordinates to \(a,b,X,Q_N\).  Because \(3^X\) is odd, the parity of
\(m_N=b+3^XQ_N\) fixes the parity of \(Q_N\) once \(b\) and the next required
bit \(x_N\) are fixed.  There are exactly four cases.

If \(b\) is even and \(x_N=0\), write \(Q_N=2Q_{N+1}\).  Then
\[
 a_N=a,\qquad b_N=b/2,\qquad X_N=X.
\]
If \(b\) is even and \(x_N=1\), write \(Q_N=1+2Q_{N+1}\).  Then
\[
 a_N=a+2^N,\qquad
 b_N=\frac{3b+3^{X+1}+1}{2},\qquad X_N=X+1.
\]
Since \(b\le3^X-1\),
\[
 b_N\le
 \frac{3(3^X-1)+3^{X+1}+1}{2}
 =3^{X+1}-1<3^{X+1}.
\]

If \(b\) is odd and \(x_N=0\), write \(Q_N=1+2Q_{N+1}\).  Then
\[
 a_N=a+2^N,\qquad
 b_N=\frac{b+3^X}{2},\qquad X_N=X,
\]
and \(b<3^X\) gives \(b_N<3^X\).  Finally, if \(b\) is odd and
\(x_N=1\), write \(Q_N=2Q_{N+1}\).  Then
\[
 a_N=a,\qquad
 b_N=\frac{3b+1}{2},\qquad X_N=X+1,
\]
and
\[
 b_N<\frac{3^{X+1}+1}{2}<3^{X+1}.
\]
In each case \(0\le a_N<2^{N+1}\), and the forced parity decomposition of
\(Q_N\) is unique.  This proves existence, uniqueness, both coefficient
bounds, and the retained quotient coordinate at stage \(N+1\)
\cite[pp.~43--44, Eq.~(4) and proof]{Everett1977}.
\end{proof}

\begin{proposition}[the exact finite parity morphism and inverse residue]
\label{prop:Everett-finite-parity-morphism}
For \(N\ge1\), define
\begin{equation}
\label{eq:Everett-PN}
 P_N:\Z/2^N\Z\longrightarrow\{0,1\}^N,
 \qquad
 P_N([m])=(T^n(m)\bmod2)_{0\le n<N}.
\end{equation}
Then \(P_N\) is a bijection.  If
\[
 c(w)=\sum_{j=0}^{N-1}x_j2^j
       3^{\sum_{i=j+1}^{N-1}x_i},
\]
its inverse is the unique residue
\begin{equation}
\label{eq:Everett-inverse-residue}
 P_N^{-1}(w)\equiv-3^{-X(w)}c(w)\pmod{2^N}.
\end{equation}
Reduction \(\Z/2^{N+1}\Z\to\Z/2^N\Z\) corresponds to deletion of the last
word digit.  On \(\Nzero\), the fibre over \(w\) is exactly
\(a(w)+2^N\Nzero\), and \eqref{eq:Everett-affine-fibre} records the quotient
that the word forgets.
\end{proposition}

\begin{proof}
Theorem~\ref{thm:Everett-finite-fibres} proves well-definedness on residue
classes, singleton fibres at residue level, surjectivity, and the stated
nonnegative fibres.  Iterating the two branch equations gives
\begin{equation}
\label{eq:Everett-word-affine-composition}
 2^NT^N(m)=3^{X(w)}m+c(w)
\end{equation}
whenever the first \(N\) parities are \(w\).  This follows by induction: the
next zero branch multiplies the existing numerator by no new factor of \(3\),
while the next one branch multiplies it by \(3\) and adds \(2^N\).
Divisibility of the right side by \(2^N\) gives
\eqref{eq:Everett-inverse-residue}; \(3^{X(w)}\) is odd and therefore a unit
modulo \(2^N\).  The same prefix read at levels \(N+1\) and \(N\) proves the
reduction identity.
\end{proof}

The source corollary uses the \(2^N\) positive representatives
\(1,\ldots,2^N\), replacing the zero residue by \(2^N\).  This implicit
representative change is valid, but it must be stated.  Numerically encoding
the target word identifies \(P_N\) exactly with the finite chronological
parity map \(\overline Q_{3,N}\) in
\eqref{eq:LDW-finite-parity-map}.  Everett proves the bijection and its
integer affine fibres.  Later sources add permutation-order, lifting,
inverse-limit, and graph-relation structure.  Everett's quotient \(Q_N\) in
the proof of Theorem~1 is not the later parity encoder \(Q_3\).

If two nonnegative integers have the same infinite parity sequence, the
finite theorem makes their difference divisible by \(2^N\) for every \(N\),
so they are equal.  This is the source's infinite injectivity conclusion.  It
does not assert that every infinite binary sequence is realized by an
ordinary nonnegative integer.

\begin{theorem}[Everett, Theorem 2: natural-density-one descent]
\label{thm:Everett-density-one-descent}
Let
\begin{equation}
\label{eq:Everett-AM}
 A(M)=\#\{1\le m\le M:\exists k\ge1,\ f_{\mathrm{Ev}}^k(m)<m\}.
\end{equation}
Then
\begin{equation}
\label{eq:Everett-density-limit}
 \lim_{M\to\infty}\frac{A(M)}M=1.
\end{equation}
Thus the source phrase ``almost every'' means ordinary natural density one
in the intervals \([1,M]\).
\end{theorem}

\begin{proof}
First take \(M=2^N\).  For a length-\(N\) word put
\begin{equation}
\label{eq:Everett-L-epsilon}
 X=\sum_{n=0}^{N-1}x_n,
 \qquad
 L=\frac{\log2}{\log(10/3)},
 \qquad
 \varepsilon=L-\frac12.
\end{equation}
Here \(\varepsilon>0\), since \(L>1/2\) is equivalent to \(4>10/3\).
Let \(H_N\) be the set of words satisfying
\begin{equation}
\label{eq:Everett-HN}
 \frac12-\varepsilon<\frac XN<\frac12+\varepsilon=L.
\end{equation}
Under uniform counting on \(\{0,1\}^N\), each coordinate is Bernoulli with
mean \(1/2\), and every prescribed pair of coordinate values has
\(2^{N-2}\) extensions.  Hence
\[
 \mathbb E X=N/2,
 \qquad
 \operatorname{Var}(X)=N/4.
\]
Chebyshev's inequality gives the exact estimate cited, but not proved, in the
source:
\begin{equation}
\label{eq:Everett-Chebyshev}
 \frac{\#H_N}{2^N}
 \ge1-\frac1{4\varepsilon^2N}.
\end{equation}
This is a finite uniform count transported by
Proposition~\ref{prop:Everett-finite-parity-morphism}; it does not assume that
one ordinary orbit has independent random steps.

Let \(D_N\) contain the words with \(X/N<L\).  Then \(H_N\subseteq D_N\).
If \(x_n=0\), then \(m_{n+1}/m_n=1/2\).  If \(x_n=1\) and \(m_n>1\), write
\(m_n=2Q+1\), where \(Q\ge1\); then
\begin{equation}
\label{eq:Everett-odd-ratio}
 \frac{m_{n+1}}{m_n}
 =\frac{3Q+2}{2Q+1}\le\frac53.
\end{equation}
Moreover,
\begin{equation}
\label{eq:Everett-word-contraction}
 \frac XN<L
 \quad\Longleftrightarrow\quad
 (1/2)^{N-X}(5/3)^X<1.
\end{equation}

Use the source corollary's representatives \(1\le m\le2^N\).  For a word in
\(D_N\), either some \(m_n=1<m\) with \(n<N\), or no such value occurs.  In
the latter case every odd factor satisfies \eqref{eq:Everett-odd-ratio}, so
\[
 m_N=m_0\prod_{n=0}^{N-1}\frac{m_{n+1}}{m_n}
 \le(1/2)^{N-X}(5/3)^Xm_0<m_0.
\]
The one excluded starting value is \(m=1\).  Therefore
\begin{equation}
\label{eq:Everett-dyadic-count}
 A(2^N)\ge\#D_N-1\ge\#H_N-1,
\end{equation}
and \eqref{eq:Everett-Chebyshev} gives
\begin{equation}
\label{eq:Everett-dyadic-density}
 \frac{A(2^N)}{2^N}
 \ge1-\frac1{4\varepsilon^2N}-2^{-N}\longrightarrow1.
\end{equation}

It remains to retain the source's interpolation, rather than claim that a
dyadic subsequence proves the full limit.  Put \(A_N=A(2^N)\) and
\begin{equation}
\label{eq:Everett-nN}
 n_N
 =2^N+\bigl((2^{N+1}-2^N)-(A_{N+1}-A_N)\bigr)
 =A_N+(2^{N+1}-A_{N+1}),
 \qquad 2^N\le n_N\le2^{N+1}.
\end{equation}
The difference \(n_N-2^N\) is exactly the number of failures in the dyadic
shell \((2^N,2^{N+1}]\).  Its bounds follow because that shell contains
\(2^N\) integers.
For \(2^N<M\le n_N\), monotonicity gives
\[
 \frac{A(M)}M\ge\frac{A_N}{n_N}.
\]
For \(M=n_N+k\), the dyadic interval has exactly
\(n_N-2^N\) failures in total, so its first
\(n_N-2^N+k\) members contain at least \(k\) successes.  Consequently
\[
 \frac{A(M)}M
 \ge\frac{A_N+k}{n_N+k}
 \ge\frac{A_N}{n_N},
\]
where the last inequality is equivalent to \(k(n_N-A_N)\ge0\).  Finally,
\begin{equation}
\label{eq:Everett-interpolation-limit}
 \frac{A_N}{n_N}
 =\frac{A_N/2^N}
 {A_N/2^N+2(1-A_{N+1}/2^{N+1})}\longrightarrow1
\end{equation}
by \eqref{eq:Everett-dyadic-density}.  The two interval estimates and
\eqref{eq:Everett-interpolation-limit} prove
\eqref{eq:Everett-density-limit}
\cite[pp.~44--45, Eqs.~(5)--(12) and Thm.~2]{Everett1977}.
\end{proof}

\begin{calculation}[an exact rational concentration witness]
\label{calc:Everett-rational-bound}
The strict inequality \(L>23/40\) is equivalent to
\[
 2^{40}3^{23}>10^{23},
\]
and direct integer subtraction gives
\[
 2^{40}3^{23}-10^{23}=3511519796081828298752>0.
\]
Thus \(\varepsilon>3/40\).  Applying the same Chebyshev proof with the
smaller rational window \(3/40\) gives the completely explicit edition
bound
\begin{equation}
\label{eq:Everett-rational-density-bound}
 \frac{A(2^N)}{2^N}
 \ge1-\frac{400}{9N}-2^{-N}.
\end{equation}
This is a proof-checking witness for the source limit, not an attribution of
the displayed constant to Everett.
\end{calculation}

\begin{proposition}[Everett--Crandall first-descent crosswalk]
\label{prop:Everett-Crandall-crosswalk}
Define
\begin{align*}
 D_T&=\{m\in\Opos:\exists k\ge1,\ T^k(m)<m\},\\
 D_C&=\{m\in\Opos:\exists j\ge1,\ C_{\mathrm{Cr}}^j(m)<m\}.
\end{align*}
Then \(D_T=D_C\) pointwise.  Moreover,
\begin{equation}
\label{eq:Everett-Crandall-relative-density}
 \lim_{M\to\infty}
 \frac{\#(D_C\cap[1,M])}{\#(\Opos\cap[1,M])}=1.
\end{equation}
This is the exact primary verification of the theorem reported by Crandall.
\end{proposition}

\begin{proof}
If \(C_{\mathrm{Cr}}^j(m)<m\), then
\(C_{\mathrm{Cr}}^j(m)=T^{A_j}(m)\) by
\eqref{eq:Crandall-three-clocks}, so \(m\in D_T\).  Conversely, suppose
\(T^k(m)<m\).  If \(T^k(m)\) is odd, it is one of the successive odd-return
states.  If it is even, put \(r=\vtwo(T^k(m))\ge1\).  The next \(r\)
shortened steps are divisions by two, and
\[
 T^{k+r}(m)=\frac{T^k(m)}{2^r}<m
\]
is odd.  It is therefore a later \(C_{\mathrm{Cr}}\)-iterate.  This proves
the pointwise equality without identifying shortened time \(k\) with
accelerated time \(j\).

Every positive even \(m\) belongs to the full shortened descent set because
\(T(m)=m/2<m\).  Put
\[
 B(M)=\#(D_C\cap[1,M]),
 \qquad
 O(M)=\#(\Opos\cap[1,M])=\lceil M/2\rceil.
\]
The exact finite-cutoff identity is therefore
\begin{equation}
\label{eq:Everett-Crandall-finite-count}
 A(M)=\lfloor M/2\rfloor+B(M).
\end{equation}
Consequently
\[
 1-\frac{B(M)}{O(M)}
 =\frac{M}{O(M)}\left(1-\frac{A(M)}M\right).
\]
The multiplier tends to \(2\).  Hence Everett's ordinary-density statement
and Crandall's relative-density statement are equivalent, and
\eqref{eq:Everett-Crandall-relative-density} follows.
\end{proof}

\begin{nonclaim}
Everett proves first descent below the starting value, not convergence to
\(1\).  The exceptional set may be infinite of density zero.  The proof gives
no uniform descent-time or total-stopping-time bound and no accelerated-time
bound.  Its finite parity surjectivity is not ordinary-integer surjectivity
onto infinite binary sequences.  It proves no statement about \(U\),
generalized \(qx+r\) maps, nontrivial cycles, or unbounded trajectories.  The
introductory verification ``up to the millions'' is historical context
without a supplied computational certificate.
\end{nonclaim}

\begin{sourceissue}[manifestation and presentation boundaries]
\label{sourceissue:Everett-manifestation}
The four raster pages are clear, but their OCR layer corrupts formulas and is
not used as mathematical evidence.  The publisher identity is fixed by DOI
and PII; the inspected content witness is a matching VOR facsimile, while a
publisher-hosted PDF hash remains unavailable.  In the article itself, the
abstract's phrase ``every integer'' must be read as every positive integer:
the fixed point \(0\) has no iterate equal to \(1\).  The body correctly uses
\(m\ge1\) for that conjecture.  The positive-representative swap and
\(Q_N\in\Nzero\) are implicit, and the probability estimate is cited rather
than proved.  The preceding reconstruction makes each step explicit.  No
theorem-level mathematical error was found.
\end{sourceissue}
